\section{System \& Speicher}

\subsection[Bestandteile]{Bestandteile $\mu$C}

\textbf{Hauptbestandteile}

\begin{itemize}
    \item CPU ($\mu$P), Speicher (ROM, RAM), Systemkomponenten (Clock).
    \item Peripherie-Elemente (GPIO, ADC, DAC, Timer, SPI, I2C, \dots)
\end{itemize}

\medbreak

\textbf{Grundelemente}

\begin{itemize}
    \item Stromquelle, POR, Taktgeber \& Boot-Code (init of internal system)
\end{itemize}

\subsection{Vorteile des Cortex-M4 NVIC}
\begin{itemize}
    \item Kein Software-Overhead für Stacking und Sprünge.
    \item Parallele Verarbeitung von Code und Daten (Harvard).
    \item Sehr geringe Jitter-Zeiten bei der Ausführung.
    \item Nur 12 Cycles für Ein- und Austritt der ESR bzw. ISR.
\end{itemize}

\vspace{0.5em}
\noindent\framebox{\parbox{0.95\columnwidth}{\centering \vspace{5pt}\textit{[Placeholder: Bild Vector-Tabelle / IRQ ISR Handler / Background/Foreground Programm]}\vspace{5pt}}}
\vspace{0.5em}

\subsection{Infos über den D-Code Systembus}
\begin{itemize}
    \item Datenzugriff für ROM (Flash schreiben).
    \item Debug-Infos aus dem ROM (Code) auslesen.
\end{itemize}

\vspace{0.5em}
\noindent\framebox{\parbox{0.95\linewidth}{\centering \vspace{5pt}\textit{[Placeholder: Speichertabelle (LA 1-28 Werte)]}\vspace{5pt}}}
\vspace{0.5em}